% preamble.tex —— 《局部 Langlands 对应的几何化》中译注本 全局导言
% 编译：xelatex ×2（有索引时 xelatex → makeindex → xelatex ×2）
\usepackage[a4paper,margin=2.6cm,footskip=1.2cm]{geometry}
\usepackage{ctex}
\usepackage{amsmath,amssymb,amsthm,mathtools}
\usepackage{mathrsfs,bbm}
\usepackage{tikz}
\usepackage{tikz-cd}
\usetikzlibrary{arrows.meta,decorations.pathmorphing,calc}
\usepackage{xcolor}
\usepackage{enumitem}
\usepackage{booktabs,array,multirow}
\usepackage{graphicx}
\usepackage{marginnote}
\usepackage[unicode,bookmarksnumbered,bookmarksopen=true]{hyperref}
\usepackage{imakeidx}
\makeindex[intoc]

% ---------- 译者体例开关：\showyizhufalse 可一键隐藏全部译注 ----------
\newif\ifshowyizhu
\showyizhutrue

% ---------- 颜色 ----------
\definecolor{yizhuGray}{gray}{0.30}
\definecolor{srcBlue}{RGB}{20,70,140}

% ---------- 译注：\yizhu{内容}{来源} ----------
% 体例：以脚注形式排，楷体，来源用 [来源：…] 标注；无外部依据写 [译者按（无外部依据）]
\newcommand{\yizhu}[2]{%
  \ifshowyizhu
    \footnote{\textcolor{yizhuGray}{\kaishu #1}\quad{\scriptsize\textcolor{srcBlue}{[来源：#2]}}}%
  \fi}
\newcommand{\yizhuw}[1]{%
  \ifshowyizhu
    \footnote{\textcolor{yizhuGray}{\kaishu #1}\quad{\scriptsize\textcolor{srcBlue}{[译者按（无外部依据）]}}}%
  \fi}

% ---------- 原文页码锚点：\srcpage{7} ----------
% 1) 版式上用 marginnote（比 marginpar 稳：跨页段、定理环境、脚注附近都不易丢）；
% 2) 同时把页码写进 <jobname>.pagemarks，任何不可见的情况都能事后核查（工具见 tools/check_pagemarks.py）。
\newif\ifsrcmark
\srcmarktrue
\newwrite\srcpageout
\immediate\openout\srcpageout=\jobname.pagemarks
\newcommand{\srcpage}[1]{%
  \ifsrcmark\immediate\write\srcpageout{#1}\fi
  \phantomsection\label{srcpage:#1}%
  \marginnote{\raggedleft\scriptsize\textcolor{yizhuGray}{p.#1}}}
% 无 margin 的精简版（用于窄页/表格内）：\pgref{N}
\newcommand{\pgref}[1]{\textsuperscript{\scriptsize\textcolor{yizhuGray}{[p.#1]}}}
% 页界落在句子中间时用这个：行内小号上标，保证不会被排版吞掉
\newcommand{\srcpageinl}[1]{%
  \ifsrcmark\immediate\write\srcpageout{#1}\fi
  \phantomsection\label{srcpage:#1}%
  \textsuperscript{\tiny\textcolor{yizhuGray}{〔p.#1〕}}}


% ---------- 编号：章用罗马数字，节/定理沿用原文「I.6.3」式编号 ----------
\renewcommand{\thechapter}{\Roman{chapter}}
\renewcommand{\thesection}{\thechapter.\arabic{section}}

% ---------- 定理环境（编号保持「章.节.序」） ----------
\theoremstyle{plain}
\newtheorem{theorem}{定理}[section]
\newtheorem{proposition}[theorem]{命题}
\newtheorem{lemma}[theorem]{引理}
\newtheorem{corollary}[theorem]{推论}
\theoremstyle{definition}
\newtheorem{definition}[theorem]{定义}
\newtheorem{example}[theorem]{例}
\newtheorem{construction}[theorem]{构造}
\newtheorem{question}[theorem]{问题}
\theoremstyle{remark}
\newtheorem{remark}[theorem]{注}
\newtheorem{conjecture}[theorem]{猜想}

% ---------- 常用记号 ----------
\newcommand{\FF}{\mathbb{F}}
\newcommand{\QQ}{\mathbb{Q}}
\newcommand{\ZZ}{\mathbb{Z}}
\newcommand{\RR}{\mathbb{R}}
\newcommand{\CC}{\mathbb{C}}
\newcommand{\PP}{\mathbb{P}}
\renewcommand{\AA}{\mathbb{A}}% 覆盖 Å
\newcommand{\GG}{\mathbb{G}}
\newcommand{\mm}{\mathfrak{m}}
\newcommand{\pp}{\mathfrak{p}}
\newcommand{\oo}{\mathcal{O}}
\newcommand{\Bun}{\mathrm{Bun}}
\newcommand{\Spa}{\mathrm{Spa}}
\newcommand{\SpaS}{\mathrm{Spa}}
\newcommand{\Div}{\mathrm{Div}}
\newcommand{\Gr}{\mathrm{Gr}}
\newcommand{\Hck}{\mathrm{Hck}}
\newcommand{\perf}{\mathrm{perf}}
\newcommand{\Adic}{\mathrm{Adic}}
\newcommand{\Crys}{\mathrm{Crys}}
\newcommand{\Isoc}{\mathrm{Isoc}}
\newcommand{\Vect}{\mathrm{Vect}}
\newcommand{\LocSys}{\mathrm{LocSys}}
\newcommand{\Rep}{\mathrm{Rep}}
\newcommand{\QCoh}{\mathrm{QCoh}}
\newcommand{\IndCoh}{\mathrm{IndCoh}}
\newcommand{\Dmod}{D\text{-}\mathrm{mod}}
\newcommand{\cD}{\mathcal{D}}
\newcommand{\cO}{\mathcal{O}}
\newcommand{\cF}{\mathcal{F}}
\newcommand{\cG}{\mathcal{G}}
\newcommand{\cH}{\mathcal{H}}
\newcommand{\cL}{\mathcal{L}}
\newcommand{\cM}{\mathcal{M}}
\newcommand{\cX}{\mathcal{X}}
\newcommand{\cY}{\mathcal{Y}}
\newcommand{\cZ}{\mathcal{Z}}
\newcommand{\wt}[1]{\widetilde{#1}}
\newcommand{\ol}[1]{\overline{#1}}
\newcommand{\ul}[1]{\underline{#1}}
\DeclareMathOperator{\Spec}{Spec}
\DeclareMathOperator{\Aut}{Aut}
\DeclareMathOperator{\Hom}{Hom}
\DeclareMathOperator{\End}{End}
\DeclareMathOperator{\Ext}{Ext}
\DeclareMathOperator{\Tor}{Tor}
\DeclareMathOperator{\id}{id}
\DeclareMathOperator{\im}{im}
\DeclareMathOperator{\coker}{coker}
\DeclareMathOperator{\rk}{rk}
\DeclareMathOperator{\Pic}{Pic}
\DeclareMathOperator{\Lie}{Lie}
\DeclareMathOperator{\Sym}{Sym}
\DeclareMathOperator{\Ind}{Ind}
\DeclareMathOperator{\Res}{Res}
\DeclareMathOperator{\Gal}{Gal}
\DeclareMathOperator{\Weil}{Weil}
\DeclareMathOperator{\Frob}{Frob}
\DeclareMathOperator{\Sht}{Sht}
\DeclareMathOperator{\Art}{Art}

% 交叉引用：\cref 风格
\newcommand{\thmref}[1]{定理~\ref{#1}}
\newcommand{\propref}[1]{命题~\ref{#1}}
\newcommand{\lemref}[1]{引理~\ref{#1}}
\newcommand{\defref}[1]{定义~\ref{#1}}
\newcommand{\remref}[1]{注~\ref{#1}}
